% Options for packages loaded elsewhere
\PassOptionsToPackage{unicode}{hyperref}
\PassOptionsToPackage{hyphens}{url}
\documentclass[11pt,a4paper]{article}
\usepackage[margin=2cm]{geometry}
\usepackage{xcolor}
\usepackage{amsmath,amssymb}
\usepackage{titlesec}
\setcounter{secnumdepth}{-\maxdimen} % remove section numbering
\usepackage{iftex}
\ifPDFTeX
  \usepackage[T1]{fontenc}
  \usepackage[utf8]{inputenc}
  \usepackage{textcomp} % provide euro and other symbols
\else % if luatex or xetex
  \usepackage{unicode-math} % this also loads fontspec
  \defaultfontfeatures{Scale=MatchLowercase}
  \defaultfontfeatures[\rmfamily]{Ligatures=TeX,Scale=1}
\fi
\usepackage{lmodern}
\ifPDFTeX\else
  % xetex/luatex font selection
\fi
% Use upquote if available, for straight quotes in verbatim environments
\IfFileExists{upquote.sty}{\usepackage{upquote}}{}
\IfFileExists{microtype.sty}{% use microtype if available
  \usepackage[]{microtype}
  \UseMicrotypeSet[protrusion]{basicmath} % disable protrusion for tt fonts
}{}
\makeatletter
\@ifundefined{KOMAClassName}{% if non-KOMA class
  \IfFileExists{parskip.sty}{%
    \usepackage{parskip}
  }{% else
    \setlength{\parindent}{0pt}
    \setlength{\parskip}{6pt plus 2pt minus 1pt}}
}{% if KOMA class
  \KOMAoptions{parskip=half}}
\makeatother
\usepackage{longtable,booktabs,array}
\newcounter{none} % for unnumbered tables
\usepackage{calc} % for calculating minipage widths
% Correct order of tables after \paragraph or \subparagraph
\usepackage{etoolbox}
\makeatletter
\patchcmd\longtable{\par}{\if@noskipsec\mbox{}\fi\par}{}{}
\makeatother
% Allow footnotes in longtable head/foot
\IfFileExists{footnotehyper.sty}{\usepackage{footnotehyper}}{\usepackage{footnote}}
\makesavenoteenv{longtable}
\setlength{\emergencystretch}{6em} % prevent overfull lines
\setlength{\parindent}{0pt}
\setlength{\parskip}{6pt plus 2pt minus 1pt}
\linespread{0.97}
\titleformat{\section}{\centering\Large\bfseries}{}{0pt}{}
\titleformat{\subsection}{\centering\large\bfseries}{}{0pt}{}
\titleformat{\subsubsection}{\centering\normalsize\bfseries}{}{0pt}{}
\titlespacing*{\section}{0pt}{1.8em}{0.7em}
\titlespacing*{\subsection}{0pt}{1.25em}{0.5em}
\titlespacing*{\subsubsection}{0pt}{0.8em}{0.3em}
\providecommand{\tightlist}{%
  \setlength{\itemsep}{0pt}\setlength{\parskip}{0pt}}
\usepackage{bookmark}
\IfFileExists{xurl.sty}{\usepackage{xurl}}{} % add URL line breaks if available
\urlstyle{same}
% fallback for those not using the hyperref driver hyperxmp:
\makeatletter
\@ifundefined{xmpquote}{\newcommand{\xmpquote}[1]{#1}}{}
\makeatother
\hypersetup{
  hidelinks,
  pdfcreator={LaTeX via pandoc}}

\author{}
\date{}

\begin{document}
\small
\sloppy

\section{Projeto de Pesquisa - Group A}\label{projeto-de-pesquisa}

\title{Integrantes: Fabio, Tobias, Vítor}


\subsection{Título}\label{tuxedtulo}

Planejamento de Visitas Domiciliares na Atenção Primária à Saúde: uma
Heurística de Horizonte Móvel para Otimização de Rotas (PHVD-APS)

\subsection{Motivação/Introdução}\label{motivauxe7uxe3ointroduuxe7uxe3o}

A Atenção Primária à Saúde (APS) é reconhecida como o eixo estruturante
dos sistemas de saúde universais. No modelo brasileiro, a Estratégia
Saúde da Família organiza equipes multiprofissionais responsáveis por
populações adscritas em territórios definidos, com ênfase no
acompanhamento longitudinal de condições crônicas. Uma das ações
centrais dessas equipes são as \textbf{visitas domiciliares}, que
garantem o contato com pacientes que não conseguem se deslocar até a
unidade de saúde e permitem identificar situações de risco, acompanhar
condições como diabetes e hipertensão e prevenir agravamentos.

Na prática, porém, o planejamento dessas visitas é realizado de forma
manual ou sem suporte computacional adequado. Isso gera acúmulo de
pendências --- visitas que deveriam ter ocorrido e não ocorreram ---,
distribuição desigual da carga de trabalho entre equipes e rotas
ineficientes que aumentam o tempo de deslocamento em detrimento do tempo
de cuidado. Quando uma visita não é realizada, o atraso se acumula para
as visitas seguintes, comprometendo o acompanhamento periódico exigido
por cada condição clínica.

O problema de planejamento de visitas domiciliares combina quatro
decisões interdependentes: \textbf{seleção de quais visitas realizar},
\textbf{escolha do dia de atendimento}, \textbf{atribuição a uma equipe}
e \textbf{ordenação da sequência de casas dentro de cada rota}. Em
horizontes de múltiplos dias úteis, a decisão tomada hoje afeta a carga
e os prazos dos dias seguintes. A reação a eventos reais --- visitas
concluídas ou não realizadas --- exige \textbf{replanejamento} que
preserve o histórico e recalcule apenas os dias futuros.

Este projeto propõe construir e avaliar um \textbf{artefato de pesquisa
reutilizável} que, a partir dos dados de uma unidade de saúde, produza
rotas para os próximos N dias de trabalho, priorize visitas pendentes,
reduza o custo de deslocamento e atualize o planejamento ao registrar o
resultado real de cada visita. O artefato inclui uma \textbf{interface
gráfica mínima} com cadastros, mapa interativo com tiles do
OpenStreetMap (via Leaflet.js ou Folium), visualização das rotas e
exportação em CSV e GPX, além de um núcleo de otimização executável em
lote para os experimentos.

\textbf{Pergunta de pesquisa:} em instâncias representativas de visitas
domiciliares e falhas de atendimento, quanto uma heurística para uma
janela móvel de N dias melhora o atendimento das pendências, o atraso
acumulado, o deslocamento e o tempo de execução em relação a regras
simples de planejamento?

\textbf{Hipótese:} uma seleção orientada por urgência e prazos futuros,
seguida de inserção de menor custo e melhoria local, reduz o atraso
acumulado após replanejamentos sem elevar excessivamente o deslocamento
ou o tempo computacional.

\subsection{Objetivos}\label{objetivos}

\subsubsection{Objetivo geral}\label{objetivo-geral}

Construir, implementar e avaliar experimentalmente uma heurística de
horizonte móvel para o planejamento de visitas domiciliares na APS,
comparando-a a baselines simples em instâncias sintéticas que
representam diferentes cenários de capacidade, distribuição geográfica e
frequência de falhas no atendimento.

\subsubsection{Objetivos específicos}\label{objetivos-especuxedficos}

\begin{enumerate}
\def\labelenumi{\arabic{enumi}.}
\tightlist
\item
  \textbf{Modelar} o problema de planejamento de visitas domiciliares
  como um problema de seleção, roteamento e atribuição de visitas a
  equipes e dias, com restrições de capacidade de jornada, prazos por
  condição clínica e replanejamento após eventos reais.
\item
  \textbf{Implementar} dois baselines diários (urgência e geográfico) e
  uma heurística principal de horizonte móvel com antecipação, além de
  melhoria local por 1.5-opt.
\item
  \textbf{Construir} um gerador de cenários sintéticos reproduzíveis, um
  validador de planos e um avaliador de métricas, de modo que todos os
  experimentos sejam replicáveis a partir de semente e configuração.
\item
  \textbf{Desenvolver} uma interface gráfica mínima com mapa interativo
  (tiles do OpenStreetMap), cadastros de posto, pacientes e equipes,
  visualização das rotas por dia e equipe, registro dos resultados das
  visitas e exportação das rotas em CSV e GPX.
\item
  \textbf{Avaliar} os métodos em cenários com capacidade folgada,
  equilibrada e insuficiente, variando distribuição geográfica, fração
  de pendências, frequência de falhas e parâmetros da função de
  pontuação, e comparar os resultados por métricas de cobertura, atraso,
  deslocamento e tempo de execução.
\item
  \textbf{Analisar} as ameaças à validade experimental, incluindo o uso
  de distância em linha reta como aproximação do deslocamento real, e
  discutir os limites de aplicação dos resultados ao contexto
  operacional do SUS.
\end{enumerate}

\subsection{Trabalhos relacionados}\label{trabalhos-relacionados}

Os trabalhos reunidos na revisão bibliográfica situam o projeto em duas
dimensões: o contexto da APS brasileira e os métodos computacionais de
otimização de rotas.

\textbf{Organização territorial e visitas domiciliares na APS.} Faria
(2020) mostra que o território da APS não é apenas uma área geográfica,
mas um espaço de responsabilidade sanitária, vínculos e necessidades de
saúde. A territorialização define quais equipes acompanham quais
famílias e em quais intervalos --- dados diretamente utilizados no
modelo proposto. Matta e Morosini (Fiocruz) e o documento CONASS (2021)
reforçam que a Estratégia Saúde da Família exige acompanhamento
longitudinal e periódico de condições crônicas, justificando o controle
de intervalos máximos entre visitas como restrição central do problema.
Rodrigues et al. (2014) identificam a falta de sistemas de informação
capazes de acompanhar pendências, intervalos e encaminhamentos como uma
das principais fragilidades da APS na coordenação do cuidado. Barros,
Aquino e Souza (2022) evidenciam a heterogeneidade entre municípios e a
relação entre privação socioeconômica e resultados de saúde, motivando
cenários experimentais com distribuições geográficas e cargas distintas.

\textbf{Tecnologias de informação na APS.} Bender et al. (2021) mostram
que, no Brasil, o uso de TIC na APS ainda é desigual e sofre com
conectividade limitada, o que justifica a exportação de rotas em
formatos portáteis (CSV, GPX) e o uso de mapas de fundo sem dependência
de servidor próprio. Pinto e Rocha (2024) descrevem o observatório
OTICS-RIO como exemplo de sistema que transforma registros cotidianos em
informação de gestão, mostrando que a rastreabilidade das decisões de
planejamento tem valor além da logística.

\textbf{Otimização de rotas para saúde comunitária.} Randriamihaja et
al. (2024) apresentam a solução mais diretamente relacionada ao projeto:
combinam dados do OpenStreetMap com um algoritmo de roteamento com
restrição de jornada (VRPTW) para distribuir visitas de agentes de saúde
em Madagáscar. Os autores mostram que distância euclidiana subestima
consideravelmente os trajetos reais e que a escolha do número de agentes
e da frequência das visitas tem impacto direto na cobertura. Em relação
ao projeto proposto, as diferenças são: (a) este projeto adota distância
em linha reta como aproximação experimental e deixa a rede viária como
extensão futura, já que a arquitetura aceita uma matriz de custos
externa; (b) este projeto incorpora explicitamente o replanejamento após
eventos reais e o controle de versões do plano; (c) este projeto compara
múltiplos métodos em vez de avaliar um único algoritmo.

A literatura de roteamento de veículos (VRP e variantes), roteamento com
lucros e planejamento de atenção domiciliar ainda precisa ser
incorporada à revisão para situar formalmente a heurística proposta e
comparar as garantias de qualidade dos métodos. Essa complementação está
prevista como parte da Etapa 5 do plano de trabalho.

\subsection{Metodologia}\label{metodologia}

A pesquisa tem abordagem \textbf{quantitativa experimental} com dados
sintéticos. Não serão utilizados dados identificáveis de pacientes;
todos os cenários são gerados computacionalmente com sementes
registradas para garantir reprodutibilidade.

\textbf{Dados e instâncias.} O gerador de cenários produz postos de
saúde, pacientes com coordenadas, condições e datas de última visita,
equipes com disponibilidade variável, polígonos de atuação, calendários
e sequências de eventos (conclusões e visitas não realizadas) com taxas
controladas. Cada cenário inclui semente, configuração completa e
sequência de eventos fixados previamente para que métodos diferentes
encontrem o mesmo evento ao programar a mesma visita.

\textbf{Cenários experimentais.} Os experimentos cobrem: capacidade
folgada, equilibrada e insuficiente; pacientes próximos e espalhados;
poucos e muitos atrasos; equipes com jornadas iguais e diferentes;
horizontes curtos e longos; falhas isoladas e sucessivas. Casos-limite
também são avaliados: zero pacientes, zero equipes, paciente fora do
polígono, condição múltipla, visita mais longa que qualquer jornada,
tentativa não realizada no último dia.

\textbf{Métodos implementados e comparados:}

\begin{enumerate}
\def\labelenumi{\arabic{enumi}.}
\tightlist
\item
  \emph{Baseline diário de urgência:} em cada dia, ordenar candidatos
  por vencimento e inserir cada um na equipe com menor custo adicional
  viável.
\item
  \emph{Baseline diário geográfico:} construir cada rota pelo vizinho
  viável mais próximo, mantendo capacidade e regras de recorrência.
\item
  \emph{Heurística principal com antecipação:} pontuar candidatos por
  atraso, peso clínico, proximidade do prazo e custo incremental de
  inserção; permitir antecipação dentro do limite A; impedir que visita
  futura de baixo peso desloque visita vencida.
\item
  \emph{Melhoria local (1.5-opt):} combinar reinserção de nós e inversão
  de subsegmentos para reduzir deslocamento dentro de cada rota sem
  alterar os dias de atendimento.
\item
  \emph{Referência exata em instâncias pequenas (opcional):} estimar a
  distância das heurísticas ao melhor resultado conhecido.
\end{enumerate}

\textbf{Avaliação.} Todos os métodos recebem o mesmo estado inicial, o
mesmo calendário, o mesmo limite de tempo e a mesma sequência de eventos
reais. As métricas são: fração de pendências efetivamente concluídas;
dias de atraso acumulados; atraso remanescente ao fim da janela; visitas
concluídas e não realizadas; deslocamento efetivo e planejado;
utilização e desequilíbrio de jornada; alterações entre versões; tempo
de geração e replanejamento; memória máxima. Resultados do plano
previsto e do simulado são separados. A comparação usa medianas e
dispersão por cenário, além de análise de dominância de Pareto para
evitar assumir um peso universal entre objetivos.

\textbf{Ferramentas e ambiente.} O núcleo de otimização é implementado
em Python, executável sem a interface para permitir testes em lote. A
interface usa Leaflet.js com tiles do OpenStreetMap como camada visual
de fundo; os traçados das rotas são segmentos em linha reta sobrepostos
a esse fundo. A matriz de deslocamento usa distância de Haversine com
velocidade fixa documentada como aproximação experimental. A
persistência salva cenário, histórico real e todas as versões do plano
em formato documentado.

\textbf{Análise de complexidade.} A análise assintótica esperada por
etapa está documentada no plano de desenvolvimento, com premissas
explícitas. Os limites serão verificados empiricamente variando n
(pacientes), m (equipes), N (dias) e o número de eventos de
replanejamento.

\subsection{Etapas e Cronograma}\label{etapas-e-cronograma}

{\def\LTcaptype{none} % do not increment counter
\begin{longtable}[]{@{}>{\raggedright\arraybackslash}p{0.12\textwidth}>{\raggedright\arraybackslash}p{0.82\textwidth}@{}}
\toprule\noalign{}
Semana & Objetivo \\
\midrule\noalign{}
\endhead
\bottomrule\noalign{}
\endlastfoot
Semana 1 & \textbf{Revisão bibliográfica e Entendimento do Domínio} ---
levantamento dos papers de APS e roteirização; fixar formato de
entrada/saída, estados de visita, calendário, geometria, restrições e
métricas; elaborar exemplo manual de N dias com falha e
replanejamento. \\
Semana 2 & \textbf{Concepção do projeto} --- redigir projeto de pesquisa
e plano de desenvolvimento; definir pergunta de pesquisa, hipótese,
modelo de dados, restrições, métricas e protocolo experimental; revisar
literatura complementar de roteamento. \\
Semana 3 & \textbf{Núcleo de otimização e interface} --- implementar
gerador de cenários, validador, estado temporal, matriz de custos,
verificador de planos, baselines de urgência e geográfico, heurística
com antecipação, 1.5-opt, mini CRUD, mapa interativo com tiles OSM,
visualização das rotas, registro de resultados e exportação em CSV e
GPX. \\
Semana 4 & \textbf{Experimentos, avaliação e escrita do artigo} ---
executar experimentos de escala, falhas e sensibilidade em todos os
métodos; calcular métricas, comparar resultados, construir tabelas e
gráficos; interpretar limites e ameaças à validade; redigir e revisar o
artigo final. \\
\end{longtable}
}

\subsection{Recursos}\label{recursos}

{\def\LTcaptype{none} % do not increment counter
\begin{longtable}[]{@{}>{\raggedright\arraybackslash}p{0.12\textwidth}>{\raggedright\arraybackslash}p{0.25\textwidth}>{\raggedright\arraybackslash}p{0.25\textwidth}>{\raggedright\arraybackslash}p{0.25\textwidth}@{}}
\toprule\noalign{}
Semana & Fábio & Tobias & Vitor \\
\midrule\noalign{}
\endhead
\bottomrule\noalign{}
\endlastfoot
Semana 1 & \textbf{Revisão bibliográfica} & \textbf{Revisão
bibliográfica} & \textbf{Revisão bibliográfica} \\
Semana 2 & \textbf{Projeto de pesquisa e plano de desenvolvimento} &
\textbf{Projeto de pesquisa e plano de desenvolvimento} &
\textbf{Projeto de pesquisa e plano de desenvolvimento} \\
Semana 3 & \textbf{Experimentação e Construção do artefato} &
\textbf{Experimentação e Construção do artefato} &
\textbf{Experimentação e Construção do artefato} \\
Semana 4 & \textbf{Análise dos resultados e redação do artigo} &
\textbf{Análise dos resultados e redação do artigo} & \textbf{Análise
dos resultados e redação do artigo} \\
\end{longtable}
}

\textbf{Recursos materiais e infraestrutura:} computadores pessoais dos
integrantes; Python 3.x com bibliotecas numpy, pandas, shapely, gpxpy e
folium/leaflet; repositório Git compartilhado para controle de versão do
código e dos cenários experimentais; nenhum dado clínico real será
utilizado.

\subsection{Referências}\label{referuxeancias}

BARROS, Rafael Damasceno de; AQUINO, Rosana; SOUZA, Luis Eugênio Portela
Fernandes. Evolução da estrutura e resultados da Atenção Primária à
Saúde no Brasil entre 2008 e 2019. \textbf{Ciência \& Saúde Coletiva},
2022. Disponível em:
\url{https://www.scielo.br/j/csc/a/rRCVJhncQt95Db9xfMxW6TF/?lang=pt}

BENDER, et al. O uso de Tecnologias de Informação e Comunicação em Saúde
na Atenção Primária à Saúde no Brasil, de 2014 a 2018. \textbf{Ciência
\& Saúde Coletiva}, 2021. Disponível em:
\url{https://www.scielo.br/j/csc/a/CFj6GmKwqyCMHTrpNPJQLXM/abstract/?lang=pt}

CONSELHO NACIONAL DE SECRETÁRIOS DE SAÚDE (CONASS). \textbf{A Atenção
Primária à Saúde no SUS: avanços e ameaças}. CONASS Documenta, n. 38.
Brasília, 2021. Organização: Eugênio Vilaça Mendes. Disponível em:
\url{https://www.conass.org.br/biblioteca/conass-documenta-38/}

FARIA, Rivaldo Mauro de. A territorialização da Atenção Básica à Saúde
do Sistema Único de Saúde do Brasil. \textbf{Ciência \& Saúde Coletiva},
v. 25, n. 11, 2020, p. 4521-4530. Disponível em:
\url{https://www.scielosp.org/article/csc/2020.v25n11/4521-4530/}

MATTA, Gustavo Corrêa; MOROSINI, Márcia Valéria Guimarães. Atenção
Primária à Saúde. In: \textbf{Dicionário da Educação Profissional em
Saúde}. Escola Politécnica de Saúde Joaquim Venâncio, Fiocruz.
Disponível em:
\url{https://scholar.google.com/citations?view_op=view_citation&hl=pt-BR&user=uFFqgI4AAAAJ&citation_for_view=uFFqgI4AAAAJ:M3ejUd6NZC8C}

PINTO, et al. Inovações na Atenção Primária em Saúde: o uso de
ferramentas de tecnologia de comunicação e informação para apoio à
gestão local. \textbf{Ciência \& Saúde Coletiva}, v. 29, n. 1, 2024.
Disponível em:
\url{https://www.scielosp.org/article/csc/2024.v29n1/e19882022/}

RANDRIAMIHAJA, et al. Combining OpenStreetMap mapping and route
optimization algorithms to inform the delivery of community health
interventions at the last mile. \textbf{PMC}, 2024. Disponível em:
\url{https://pmc.ncbi.nlm.nih.gov/articles/PMC11542841/}

RODRIGUES, Ludmila Barbosa Bandeira et al. A atenção primária à saúde na
coordenação das redes de atenção: uma revisão integrativa.
\textbf{Ciência \& Saúde Coletiva}, v. 19, n. 2, 2014, p. 343-352.
Disponível em:
\url{https://www.scielo.br/j/csc/a/nBKRxhLTPkdp489zfNGhKnt/abstract/?lang=pt}

\end{document}
